\section{Instruction Finetuning}

\subsection{Método básico}

Se recopilan pares (instruction, output) y se continúa con un ajuste fino supervisado sobre el modelo preentrenado.

\begin{importantbox}{Efecto del Instruction Finetuning}
Después de aplicar instruction finetuning en más de 1,800 tareas, FLAN-T5 mejora de forma notable su desempeño en tareas no vistas. Los modelos más grandes obtienen más beneficio del finetuning.
\end{importantbox}

\subsection{Estrategia de datos y alineación}

- En el nivel de los datos, primero se aplican prompt rewriting y canonicalization para reducir la redundancia semántica; después se etiquetan la mejor respuesta y las respuestas alternativas.
- Se simulan diálogos de varios turnos para aumentar la dependencia del contexto; a los ejemplos negativos se les agregan tipos de error diversos para mejorar la robustness.
- Al recopilar las instrucciones, se conserva también la user intent metadata (tipo de tarea, complejidad, preferencias) para que después el reward model pueda hacer un modelado estratificado.

\begin{knowledgebox}{Hacer interpretables los datos de instruction}
Se agrega metadata a cada instruction para que, durante el entrenamiento, pueda elegirse de manera dinámica la reward signal según la intención (por ejemplo, \textit{creative vs. factual}), lo que reduce el error de generalización causado por una sola loss.
\end{knowledgebox}

\begin{importantbox}{La calidad de los datos importa más que el tamaño del modelo}
Aunque los modelos más grandes tienen capacidades subyacentes más sólidas, la calidad de los datos del ajuste fino con instructions determina la confiabilidad del asistente final. Los campos adversariales, la diversidad de instrucciones y la revisión humana son garantías básicas.
\end{importantbox}

\subsection{Limitaciones del Instruction Finetuning}

\begin{enumerate}
    \item \textbf{La generación abierta no tiene una respuesta estándar}: tareas como la escritura creativa no pueden supervisarse con una única salida correcta.
    \item \textbf{La penalización a nivel de Token no distingue la gravedad del error}: la pérdida del modelado de lenguaje trata por igual todos los errores de token.
    \item \textbf{Las respuestas generadas por personas no son perfectas}: los datos etiquetados tienen ruido.
\end{enumerate}

\subsection{Resumen de la sección}

El instruction finetuning es el primer paso para pasar de un LM a un asistente. Debe acompañarse de una estrategia de datos y de una separación de preferencias para proporcionar una política inicial estable a RLHF o DPO.

